\begin{multicols}{3}
\byline{Споделено}{Мартина Миразчиева, Маргарита Панайотова, Х „Д“ клас}

\noindent \begin{window}[2,r, \includegraphics[width=2.1in]{./Todorova/todorova.jpg},] \end{window}

На 10.03.2015 година се срещнахме с  педагогическия съветник на Немска езикова 
гимназия Гергана Тодорова. Тя работи в гимназията от 15 септември 2014 година. 
Атмосферата в гимназията много й допада и това е една от главните причини да 
идва на работа с голямо желание. Учениците, които идват при нея, най-често 
търсят душевен комфорт, търсят спокойно и уютно място, където да намерят решение 
на затрудненията си.  Проблемите им са от всякакъв характер – проблем с 
взаимоотношенията със съучениците и учителите, но най-често срещаните са 
характерни за етапа на личностното развитие проблеми. Когато един ученик отиде 
при нея, той се надява да намери защита, сигурност и разчита на грижите и 
желанието, с което педагогическият съветник ще се отнесе към неговия проблем. 
Едно от главните неща, които не й харесват в училище, е неспазването на 
правилата. Не й харесва също неуважението от страна на учениците към учителската 
професия и намира за много притеснителна демотивацията на някои ученици към 
усвояването на знания извън строго рамкирания учебен материал. При госпожа 
Тодорова не идват само ученици, а и учители, които искат да разрешат 
недоразуменията, възникнали в клас по време на учебния процес. Смята, че 
психологът не бива да е прекалено възрастен. Все пак общува и работи с млади 
хора и възрастовата разлика не бива да е много голяма. На въпроса какво би 
променила в гимназията, госпожа Тодорова трудно ни отговори. Все пак би искала 
да подобри работата в екип. Останала е с впечатление, че по-голямата част от 
учениците са индивидуалисти и не са свикнали да работят в екип. Когато за финал 
я попитахме какъв съвет би дала на учениците, тя отвърна : „ Учете повече, в 
колкото се може повече сфери, защото никое знание не е излишно“. Разговорът ни 
беше увлекателен и интересен. Научихме много за нейната личност и за работата й.

\closearticle
\end{multicols}
